\section{Methodology}
\label{sec:method}

\subsection{Additive Prompt Levels}
\label{sec:levels}

For each task we define seven prompt templates $T_1 \subset T_2 \subset
\cdots \subset T_7$, where $\subset$ denotes strict token superset.
Each level appends one layer of elaboration, a distinct prompt
\emph{mechanism}, to all previous content; no layer modifies earlier text.
Table~\ref{tab:levels} shows the layer structure for all six tasks; the
complete prompt text at every level, with worked walkthroughs, is in the
supplementary material.

\begin{table}[t]
\centering
\caption{Additive prompt layers by task. Each row adds content to all
  previous rows. Layers represent qualitatively different prompt
  mechanisms, not homogeneous units of length.}
\label{tab:levels}
\small
\resizebox{\textwidth}{!}{
\begin{tabular}{cllllllll}
\toprule
Lv & Mechanism & Classification & Prod.\ Extraction & QA & Math Reasoning & Summarisation & Instr.\ Following \\
\midrule
1 & Bare input   & Raw text only & Product text & Question only & Problem statement & Article text & Topic prompt \\
2 & Task label   & ``Classify sentiment'' & List field names & ``Answer the question'' & ``Solve the problem'' & ``Summarise this'' & Task + constraints \\
3 & Format spec  & ``Pos/Neg/Neu'' & JSON format & ``Answer concisely'' & ``Give a number'' & Length target & Format specification \\
4 & Definitions  & Label definitions & Field definitions & Domain context & Show-work instr. & Content guidance & Constraint definitions \\
5 & Persona      & Expert analyst & Expert extractor & Domain expert & Math tutor & Journalist & Writing expert \\
6 & Guidelines   & Mixed/ironic rules & Edge cases & Uncertainty & Step-by-step rules & Quality rules & Formatting rules \\
7 & Example      & Worked example & Worked example & Worked example & Worked example & Worked example & Worked example \\
\bottomrule
\end{tabular}
}
\end{table}

The additive design is critical: because level $k+1$ contains all tokens of
level $k$ plus exactly one new mechanism, any quality difference between
adjacent levels is attributable solely to that mechanism.
No layer is filler; every added layer is meaningful instruction that could
plausibly improve performance.

\subsection{Tasks}
\label{sec:tasks}

We study six tasks spanning a spectrum of output structure:

\begin{itemize}
  \item \textbf{Classification}: sentiment classification (positive / negative
    / neutral) over 20 product reviews and social media posts.
    Heuristic evaluator: correct label present in response.
  \item \textbf{Product extraction}: extract name, price, brand, category
    from product listings in JSON format.
    20 examples with varying difficulty (buried prices, implicit brand names).
    Heuristic: fraction of four fields correctly matched.
  \item \textbf{Question answering (QA)}: factoid questions across geography,
    science, and history. 20 examples with ground-truth strings.
    Heuristic: word-overlap recall (ROUGE-1).
  \item \textbf{Math reasoning}: 2--3 step arithmetic, rate, and geometry
    problems. 20 examples (10 easy, 10 medium).
    Heuristic: exact numerical match on the last number in the response.
  \item \textbf{Summarisation}: abstractive summarisation of news passages.
    20 examples with reference summaries.
    Heuristic: length score + ROUGE-1 recall.
  \item \textbf{Instruction following}: responses must satisfy explicit
    structural constraints (bullet points, numbered list, single word).
    20 examples, each with 1--2 constraints.
    Heuristic: fraction of constraints satisfied.
  \item \textbf{Named Entity Recognition (NER, held-out validation):}
    extract person names, organisations, and locations from news-style
    sentences. 20 curated examples spanning easy (single entity type) to
    hard (multiple overlapping entities). Heuristic: per-type F1 averaged
    across entity categories.
    This task was selected \emph{after} the main study to test the
    schema-gating mechanism out of sample. We predicted, before running, that
    NER would be \textbf{schema-gated}, with its quality gain concentrated at
    the single output-format layer (L3), because the model cannot infer the
    typed JSON structure from a bare sentence.
\end{itemize}

To assess robustness to dataset size, we replicate the experiment on
200 examples for classification (SST-2; \citealt{socher2013recursive})
and QA (SQuAD; \citealt{rajpurkar2016squad}), preserving the same prompt
templates and evaluation procedure.
These two tasks represent opposite poles, schema-gated (strong saturation)
vs.\ knowledge-gated (no saturation); confirming both the positive and
negative findings at larger scale on standard benchmarks bounds the
intermediate tasks.
Product extraction receives independent confirmation through the ablation
(2,986 runs with new orderings and examples) and padding control
(1,470 runs) rather than a separate replication.

To test whether ``insensitive'' classifications reflect genuinely easy
examples rather than true prompt insensitivity, we run a hard-example
stress test: 20 multi-step GSM8K problems~\citep{cobbe2021gsm8k} for
math reasoning, 20 multi-hop HotpotQA questions~\citep{yang2018hotpotqa}
for QA, and 20 multi-constraint instructions for instruction following,
evaluated on three models across all seven levels (1,260 runs).

\subsection{Models}
\label{sec:models}

We evaluate eight LLMs across four providers, spanning 1B to 70B disclosed
parameters plus four API-tier models with undisclosed sizes
(Table~\ref{tab:models}).
All models are accessed via production APIs with temperature set to 0
(greedy decoding) for deterministic outputs.
Each response is capped at 512 output tokens.
The seven primary models run all six tasks; Llama-3.2-1B is evaluated on
classification and product extraction only, the two tasks where saturation
is present and the parameter-scale interaction is testable, to extend the
size range below 8B without inflating the experiment matrix.

\begin{table}[t]
\centering
\caption{Models evaluated. Parameter counts are listed where publicly
  disclosed; ``undiscl.''\ indicates the provider has not released this
  information. $\dagger$Evaluated on classification and product extraction
  only. $\ddagger$Deprecated between study phases; supplementary
  experiments (NER, hard examples, stability) use gemini-2.5-flash and
  exclude kimi-k2.}
\label{tab:models}
\small
\begin{tabular}{lllr}
\toprule
Model (alias) & Provider & API & Parameters \\
\midrule
llama-3.2-1b-instruct$^\dagger$ & Meta & Groq     & 1B        \\
llama-3.1-8b-instant    & Meta       & Groq      & 8B        \\
qwen/qwen3-32b          & Alibaba    & Groq      & 32B       \\
llama-3.3-70b-versatile & Meta       & Groq      & 70B       \\
gemini-2.0-flash$^\ddagger$ & Google  & Gemini    & undiscl.  \\
moonshotai/kimi-k2$^\ddagger$ & Moonshot & Groq   & undiscl.  \\
gpt-4o-mini             & OpenAI     & OpenAI    & undiscl.  \\
claude-haiku-4-5        & Anthropic  & Anthropic & undiscl.  \\
\bottomrule
\end{tabular}
\end{table}

\subsection{Evaluation}
\label{sec:eval}

We deliberately avoid resting our conclusions on a single LLM evaluator.
For every task that admits objective ground truth we use a \emph{deterministic}
metric as the \textbf{primary} measure; the LLM judge is primary only where no
clean heuristic exists, and elsewhere serves as a cross-check.

\paragraph{Primary metric: deterministic ground truth.}
Four tasks have unambiguous ground truth and are scored by exact, reproducible
rules with no model in the loop: \emph{classification} (label match),
\emph{product extraction} (fraction of the four fields exactly matched against
gold JSON), \emph{QA} (ROUGE-1 recall / answer-string containment), and
\emph{math} (exact numerical match on the final answer).
The held-out NER task uses per-type entity F1.
Because these metrics cannot be confounded by judge calibration, they are the
basis for all saturation claims on these tasks.

\paragraph{Judge metric (where no ground truth exists).}
\emph{Summarisation} and \emph{instruction following} lack a single correct
output, so following \citet{zheng2023judging} we score them with
\texttt{gpt-4o-mini} on four dimensions (correctness, completeness, reasoning,
conciseness), each rated 1--5, aggregated as
$q = \frac{\sum_i s_i}{20} \in [0,1]$, with task-specific rubrics.
We also record judge scores for the four ground-truth tasks so the two metrics
can be compared head-to-head (Section~\ref{sec:result-judge}).

\paragraph{Independent second judge.}
As a further reliability check, all 3,915 judge-scored main-study responses
are re-evaluated by \texttt{gemini-2.0-flash} using the identical rubric,
yielding inter-judge agreement statistics (Section~\ref{sec:result-judge}).

\paragraph{Why ground-truth-primary.}
Across the full dataset the judge and the deterministic metric correlate
strongly overall ($r = 0.986$), but they diverge informatively on individual
tasks, most importantly product extraction, where the judge awards partial
credit to malformed output and thereby both masks the format step and spreads
apparent gain onto the worked example.
Where the two disagree we treat the deterministic metric as the arbiter and
report the disagreement explicitly.
This design converts the ``LLM-judge reliability'' question from a threat into
evidence.

\paragraph{Aggregation.}
For each (model, task, level) triplet, we average quality across the 20
examples, yielding 7 data points per (model, task) curve.

\subsection{Curve Fitting and Saturation Points}
\label{sec:fitting}

For each (model, task) pair we fit two functional forms to the 7 level means:
\begin{equation}
  q_{\log}(x) = a \log(bx) + c, \qquad
  q_{\text{sig}}(x) = \frac{L}{1 + e^{-k(x - x_0)}} + c
  \label{eq:curves}
\end{equation}
where $x$ is the mean prompt token count at each level.
We select the form with lower RMSE.
The \textbf{saturation point} $T^*$ is defined as the smallest $x$ such that:
\begin{equation}
  q(x) \geq 0.95 \times \max_{x} q(x)
  \label{eq:satpoint}
\end{equation}

\paragraph{Null model F-test.}
To test whether a fitted curve is significantly better than a flat line
($q(x) = \bar{q}$), we compute:
\begin{equation}
  F = \frac{(SS_{\text{null}} - SS_{\text{fit}}) / (k - 1)}{SS_{\text{fit}} / (n - k)}
  \label{eq:ftest}
\end{equation}
where $k \in \{3, 4\}$ is the number of curve parameters, $n = 7$ is the
number of level means, and $\alpha = 0.05$.

\paragraph{Bootstrap confidence intervals.}
We estimate 95\% CIs on $T^*$ via 1,000-iteration bootstrap.
On each iteration we resample the 20 example indices \emph{with replacement}
and apply the same resample across all seven levels, preserving paired
structure.

\paragraph{Threshold sensitivity.}
To assess robustness of saturation estimates, we recompute $T^*$ at four
threshold levels (85\%, 90\%, 95\%, 99\% of the fitted asymptote) and
using a threshold-free second-derivative knee estimator.

\subsection{Marginal Contribution Analysis}
\label{sec:marginal}

For each task, we compute the quality delta between adjacent level pairs
($\Delta_{k \to k+1} = \bar{q}_{k+1} - \bar{q}_k$), averaged across all
models.
This decomposes the total quality gain into per-mechanism contributions,
revealing which prompt layers drive improvement and which are redundant.
We compute it on the primary (deterministic) metric for the ground-truth
tasks; the per-layer profile is what assigns a task to a class (schema-,
difficulty-, or knowledge-gated) and, crucially, detects step-shaped
concentration that curve fitting misses.

\subsection{Randomised Layer-Ordering Ablation}
\label{sec:ablation}

To disentangle information content from token count, we conduct a randomised
ablation.
For classification and product extraction, the tasks with significant
saturation in the core study, we shuffle the introduction order of all
six prompt layers (layers 2--7; layer 1 is always the bare input),
generating 5 random permutations per task.
Each permutation is run on 7 curated examples per task across 6 models
(2,986 total runs; 2,450 analysed after excluding kimi-k2 due to API
deprecation between study phases).

If saturation is driven by token count, shuffling should not affect the
quality-versus-token curve.
If saturation is driven by information content, the saturation \emph{point}
(in tokens) should shift with ordering, because high-value layers can now
appear early, while the saturation \emph{shape} (quality rising then
plateauing) should persist.

\subsection{Irrelevant-Padding Control}
\label{sec:padding}

As a direct test of the token-count null hypothesis, we replace real prompt
content with irrelevant filler matched to the same token counts.
Three padding types are used: factual trivia unrelated to the task
(geography, history, science), repeated neutral phrases, and random
English words in shuffled order.
The bare task instruction is preserved as the only meaningful content.
This yields 1,470 runs (5 models $\times$ 2 tasks $\times$ 3 padding types
$\times$ 7 levels $\times$ 7 examples).
The 7 examples per task are drawn from the main study's curated set, not a
separate sample.

If tokens themselves drive quality, padded prompts should show the same
saturation curve as real prompts at matched token counts.

\subsection{Experimental Scale}
\label{sec:scale}

Table~\ref{tab:scale} summarises the total experimental scope.

\begin{table}[t]
\centering
\caption{Experimental scale across all study components.}
\label{tab:scale}
\small
\begin{tabular}{lr}
\toprule
Component & Evaluations \\
\midrule
Main study (7 models $\times$ 6 tasks $\times$ 7 levels $\times$ 20 examples) & 5,875 \\
200-example replication (classification + QA) & $\approx$2,800 \\
Randomised layer-ordering ablation & 2,986 \\
Irrelevant-padding control & 1,470 \\
Independent second judge (gemini-2.0-flash) & 3,915 \\
Sub-8B experiment (Llama-3.2-1B) & 332 \\
Held-out NER validation (5 models $\times$ 7 levels $\times$ 20 examples) & 700 \\
Hard-example stress test (3 models $\times$ 3 tasks $\times$ 7 levels $\times$ 20 examples) & 1,260 \\
Sampling stability (3 models $\times$ 3 tasks $\times$ 7 levels $\times$ 20 examples $\times$ 2 repeats) & 2,520 \\
\midrule
\textbf{Total} & \textbf{21,858} \\
\bottomrule
\end{tabular}
\end{table}
